% TOP_PROBLEM: {"id": 30007341, "title": "Strong diffeomorphic approximation of three-dimensional homeomorphisms above exponent two", "category_id": 9, "category": {"id": 9, "name": "pde", "display_name": "Partial Differential Equations", "description": "PDEs and their applications in physics and geometry.", "slug": "pde", "order_index": 9, "created_at": "2026-07-31T15:26:25.671Z"}, "status": "open"}
% FIRST_POSTED: 2026-10-06
% ATTEMPT_STATUS: solved
% OPENAI_CLAIM: {"schema_version":1,"source_repo":"https://github.com/openai/math","source_commit":"adc7f1241b42e322a6451854ab7e4b4c146bf78a","release_date":"2026-10-06","match":"full","resolution":"proved","independently_reviewed":false,"second_pass":"agreed","families":[{"family":"368","title":"The three-dimensional Ball–Evans approximation problem","subject":"Partial differential equations","manuscripts":[{"title":"Strong diffeomorphic approximation in three dimensions for p>2","pdf_path":"preprints/Strong-diffeomorphic-approximation-in-three-dimensions-for-p-gt-2-September-24-2026/main.pdf","readme_path":"preprints/Strong-diffeomorphic-approximation-in-three-dimensions-for-p-gt-2-September-24-2026/README.md","date":"2026-09-24","theorem_ref":"main:theorem (sections/01-introduction.tex:25)","statement_sources":[{"path":"preprints/Strong-diffeomorphic-approximation-in-three-dimensions-for-p-gt-2-September-24-2026/build/sections/01-introduction.tex","line":25,"label":"main:theorem"}]}],"lean":null,"reasoning_trace":null}]}
\documentclass[11pt]{article}
\usepackage[margin=25mm]{geometry}
\usepackage{fontspec}
\setmainfont{Latin Modern Roman}
\usepackage{amsmath,amssymb,mathtools}
\usepackage{unicode-math}
\setmathfont{Latin Modern Math}
\usepackage{enumitem}
\usepackage{xurl,hyperref,bookmark}
\hypersetup{colorlinks=true,urlcolor=blue}
\setcounter{secnumdepth}{0}
\setlength{\parindent}{0pt}
\setlength{\parskip}{5pt}
\setlength{\emergencystretch}{3em}
\begin{document}
\section*{Strong diffeomorphic approximation of three-dimensional homeomorphisms above exponent two}

\subsection{OpenAI's claim}
\textit{Ball--Evans approximation above two}: For every finite $p>2$, the companion claims strong $W^{1,p}$ approximation of bounded-domain homeomorphisms in three dimensions by smooth onto diffeomorphisms. No inverse Sobolev regularity or positive Jacobian hypothesis is imposed.

\subsection{Mathematical statement of the claimed result}
\paragraph{Mathematical claim 1.} Let $\Omega,\Lambda\subset\mathbb R^3$ be nonempty bounded connected open sets and let $2<p<\infty$. For every homeomorphism $f:\Omega\to\Lambda$ in $W^{1,p}(\Omega;\mathbb R^3)$, the claim provides smooth diffeomorphisms $f_j:\Omega\to\Lambda$, each in the same Sobolev class, such that \[\int_\Omega(|f_j-f|^p+|Df_j-Df|^p)\,dx\longrightarrow0.\] The source and target are preserved exactly; no boundary regularity, prescribed boundary values or convergence of inverse derivatives is assumed.

\subsection{Outline of the argument}
\textit{Ball--Evans approximation above two}: Invertible jets admit exact affine replacements, while rank-one and rank-two regions are retained rather than discarded as volume-null. Two scalar labels control mean gradients through boundary integrals, and near-sharp norm estimates give strong convergence by uniform convexity. Locally bi-Lipschitz replacement and the included smoothing construction then assemble an onto global approximation.

\subsection{Scope relative to this problem}
\textit{Ball--Evans approximation above two}: The Sobolev exponent is any finite $p>2$, and the approximating maps are smooth diffeomorphisms onto the original target domain. The requested norm is the forward $W^{1,p}$ norm; no control of the inverse maps or of an elastic determinant penalty is claimed.

With the hypotheses and conventions above, the claimed result addresses the whole mathematical assertion on this page.

\subsection{Formal verification}
\textit{Ball--Evans approximation above two}: The pinned catalogue lists no Lean formalization for this result.

This site has not run the Lean development. These links describe OpenAI’s selected formal statements and should be read with their stated scope.

\subsection{Status caveat}
OpenAI reports that an internal model produced the manuscripts. Their arguments have not been independently verified as proofs by this site. Any portion outside the documented formal scope may contain mistakes, and the catalogue label reports the announced mathematical claim.

\subsection{Sources}
\begin{itemize}
\item OpenAI, \emph{Strong diffeomorphic approximation in three dimensions for p>2}, 2026-09-24; citation key \texttt{OAI:Strong-diffeomorphic-approximation-in-three-dimensions-for-p-gt-2-September-24-2026}. \href{https://github.com/openai/math/blob/adc7f1241b42e322a6451854ab7e4b4c146bf78a/preprints/Strong-diffeomorphic-approximation-in-three-dimensions-for-p-gt-2-September-24-2026/main.pdf}{PDF (GitHub)}; \href{https://github.com/openai/math/raw/adc7f1241b42e322a6451854ab7e4b4c146bf78a/preprints/Strong-diffeomorphic-approximation-in-three-dimensions-for-p-gt-2-September-24-2026/main.pdf}{PDF (direct)}; \href{https://github.com/openai/math/blob/adc7f1241b42e322a6451854ab7e4b4c146bf78a/preprints/Strong-diffeomorphic-approximation-in-three-dimensions-for-p-gt-2-September-24-2026/README.md}{Manuscript page and citation}; \href{https://github.com/openai/math/blob/main/preprints/Strong-diffeomorphic-approximation-in-three-dimensions-for-p-gt-2-September-24-2026/main.pdf}{Latest version}; Statement source: \href{https://github.com/openai/math/blob/adc7f1241b42e322a6451854ab7e4b4c146bf78a/preprints/Strong-diffeomorphic-approximation-in-three-dimensions-for-p-gt-2-September-24-2026/build/sections/01-introduction.tex\#L25}{main:theorem}.
\item \href{https://github.com/openai/math/blob/adc7f1241b42e322a6451854ab7e4b4c146bf78a/CONTENTS.md}{OpenAI catalogue entry}.
\end{itemize}
\end{document}
